
\section{Deterministic Problem Formulation}
By applying the previous results, we can formulate \cref{eq:problem_statement} as a nonconvex deterministic optimization problem.
\subsection{Closed-Loop Square Root Covariance Propagation}
Since the state is controlled by the affine feedback control law, the closed-loop dynamics are given for $k \in [N{-}1]$ as
\begin{equation}
	x_{k+1} = A_k x_k + B_k (v_k + K_k (x_k - \E{x_k})) + G_k w_k.
\end{equation}
It is straightforward to derive the mean/covariance dynamics:
\begin{align} \label{eq:mean_dynamics}
	\mu_{k+1} &= A_k \mu_k + B_k v_k , \quad \mu_k := \E{x_k}\\
\label{eq:full_covariance_propagation}
	P_{k+1}& = (A_k + B_k K_k) P_k (A_k + B_k K_k)^\top + G_k G_k^\top.
\end{align}
Letting $S_k := P_k^{1/2}$, define the following matrix:
\begin{equation} \label{eq:X_definition}
	X_{k+1} := \begin{bmatrix} (A_k + B_k K_k) S_k & G_k \end{bmatrix}.
\end{equation}
\cref{lemma:rank_M_M_top_rank_M,lemma:positive_definite_iff_rank_n} are used to show \cref{lemma:rank_X_n_iff_P_succ_0}. They are standard results so we omit the  proofs.
\begin{lemma} \label{lemma:rank_M_M_top_rank_M}
	For any $M \in \mathbb{R}^{p \times q}$, $\rank(M M^\top) = \rank(M).$
	% \begin{equation}
	% 	\rank(M M^\top) = \rank(M).
	% \end{equation}
\end{lemma}
\begin{lemma} \label{lemma:positive_definite_iff_rank_n}
	A symmetric positive semidefinite matrix $P$ of size $n \times n$ is strictly PD if and only if $\rank(P) = n$.
\end{lemma}
\begin{lemma} \label{lemma:rank_X_n_iff_P_succ_0}
	Let $P_{k+1}$, $X_{k+1}$ be given by \cref{eq:full_covariance_propagation}, \cref{eq:X_definition}. Then,
	\begin{equation}
		\rank(X_{k+1}) = n \Leftrightarrow P_{k+1} \succ 0.
	\end{equation}
\end{lemma}
\begin{proof}
	We use \cref{lemma:rank_M_M_top_rank_M,lemma:positive_definite_iff_rank_n} throughout the proof. \\
	($\Rightarrow$) 
	$\rank(X_{k+1}) {=} n \Rightarrow \rank(X_{k+1} X_{k+1}^\top) = n \Rightarrow P_{k+1} \succ 0$.\\
	($\Leftarrow$)
	Define $\widetilde{P}_{k+1} := (A_k + B_k K_k) P_k (A_k + B_k K_k)^\top$.
	\begin{align*}
	P_{k+1} \succ 0 &\Rightarrow \widetilde{P}_{k+1} \succ 0 \lor G_k G_k^\top \succ 0 \\
	&\Rightarrow \rank(\widetilde{P}_{k+1}) = n \lor \rank(G_k G_k^\top) = n \\
	&\Rightarrow \rank((A_k + B_k K_k) S_k) = n \lor \rank(G_k) = n \\
	&\Rightarrow \rank(X_{k+1}) = n.
	\end{align*}
	The last implication is since the column space of $X_{k+1}$ includes the column space of its submatrices.
\end{proof}

\begin{lemma}[Closed-Loop Square Root Covariance Propagation]
	\label{lemma:sqrt_propagation}
	The covariance propagation \cref{eq:full_covariance_propagation} is equivalently expressed in Cholesky factor form using $S_k = P_k^{1/2} \in \mathbb{L}_+^n$ as:
	\begin{equation} \label{eq:square_root_covariance_propagation_nonlinear_feedback_control}
		S_{k+1} = \mathrm{qr}(X_{k+1}^\top)^\top
	\end{equation}
	with $X_{k+1}$ defined as in \cref{eq:X_definition}.
\end{lemma}

\begin{proof}
The covariance propagation can be expressed as:
\begin{equation} \label{eq:covariance_propagation}
	S_{k+1} S_{k+1}^\top = X_{k+1} X_{k+1}^\top.
\end{equation}
Note that we want $S_{k+1} \in \mathbb{L}_+^n$, while $X_{k+1} \in \mathbb{R}^{n \times (n+m)}$.
For a matrix $Q_{k} \in \R^{(n+m) \times n}$ such that $Q_{k}^\top Q_k= I_n$,
\begin{subequations}
\begin{align}
     X_{k+1}^\top &= Q_{k} S_{k+1}^\top  \label{eq:Q_S_X_relation_1}\\
	\Rightarrow X_{k+1} X_{k+1}^\top &= S_{k+1} Q_{k}^\top Q_{k} S_{k+1}^\top = S_{k+1} S_{k+1}^\top.
\end{align}%
\label{eq:Q_S_X_relation}
\end{subequations}%
Since $S_{k+1} \in \mathbb{L}_+^n$, its transpose is upper triangular, so \cref{eq:Q_S_X_relation_1} is a QR decomposition of $X_{k+1}^\top$. 
Combining \cref{lemma:rank_X_n_iff_P_succ_0} and \cref{assump:positive_definite_state}, $X_{k+1}$ is full rank.
Then, from \cref{lemma:existence_and_uniqueness_of_economy_size_qr_decomposition}, there exists a unique set of matrices $(Q_{k}, S_{k+1})$ that satisfy the equation.
Thus, this QR decomposition is unique. 
\end{proof}
\cref{lemma:sqrt_propagation} ensures that $S_{k+1}$ obtained from \cref{eq:square_root_covariance_propagation_nonlinear_feedback_control} provides a unique matrix in $\mathbb{L}_n^+$ such that \cref{eq:full_covariance_propagation} is satisfied.

Since both $S_k$ and $K_k$ are unknowns, \cref{eq:square_root_covariance_propagation_nonlinear_feedback_control}, via \cref{eq:X_definition}, has a bilinear term \( K_k S_k \); for the sake of convexity, we perform a change of variables and define $L_k = K_k S_k \in \mathbb{R}^{m \times n}$. Then, $K_k$ can be uniquely recovered as $K_k = L_k S_k^{-1}$. The inverse exists since $S_k$ is the Cholesky factor of \( P_k \), which is PD by \cref{assump:positive_definite_state}. Note that $L_k L_k^\top = K_k S_k S_k^\top K_k^\top = K_k P_k K_k^\top = \cov{u_k}$. With this change of variables, we get
\begin{equation} \label{eq:square_root_covariance_propagation}
	S_{k+1} = \mathrm{qr}(X_{k+1}^\top)^\top, \ X_{k+1} = \begin{bmatrix} A_k S_k + B_k L_k & G_k \end{bmatrix}
\end{equation}
with its first-order approximation given by:
\begin{equation} \label{eq:square_root_covariance_propagation_linearized}
	\hspace{-1pt}S_{k+1} \approx \mathrm{qr}(\bar{X}_{k+1}^\top)^\top + \dd R( \Delta X_{k+1}^\top , \bar{Q}_{k+1}, \bar{R}_{k+1})^\top
\end{equation}
with \( \bar{X}_{k+1} := \begin{bmatrix} A_k \bar{S}_k + B_k \bar{L}_k & G_k \end{bmatrix} \), \(\bar{S}_k\), and \(\bar{L}_k\) being the reference values, and \(\bar{Q}_{k+1} \bar{R}_{k+1} = \bar{X}_{k+1}\) being the QR decomposition of \(\bar{X}_{k+1}\). 
Define $\Delta X_{k} := X_k - \bar{X}_k$.

\subsection{Objective Function}
After substituting \( u_k = K_k x_k = L_k S_k^{-1} x_k \), and performing some standard manipulations, the objective function can be expressed as deterministic convex functions of \( L_k \) and \( S_k \):

{
\setlength{\abovedisplayskip}{0pt} % Remove extra space above
\setlength{\belowdisplayskip}{0pt} % Remove extra space below
\small
\begin{equation} \label{eq:objective_deterministic}
	\hspace{-3pt}
	J = \begin{cases}
	\sum_{k=0}^{N-1} \mu_k^\top Q_k \mu_k + v_k^\top R_k v_k & \\
	\quad + \mathrm{trace} (Q_k S_k S_k^\top + R_k L_k L_k^\top) & \text{(EoQ)} \\
	\sum_{k=0}^{N-1} \norm{W_k^x x_k}_2 + \sqrt{\mathcal{Q}_{\chi^2_{w_x}}(p_J)} \ \norm{W_k^x S_k}_2 & \\
	\quad + \norm{W_k^u v_k}_2 + \sqrt{\mathcal{Q}_{\chi^2_{w_u}}(p_J)} \ \norm{W_k^u L_k}_2 & \text{(QoN)}.
	\end{cases}
\end{equation}
\setlength{\abovedisplayskip}{0pt} % Remove extra space above
\setlength{\belowdisplayskip}{0pt} % Remove extra space below
}

The QoN expression is a convex upper bound based on the triangle inequality \cite{oguriChanceConstrainedControlSafe2024a}; \( \mathcal{Q}_{\chi^2_d}(\cdot) \) is the inverse of the chi-squared cumulative distribution function with \( d \) degrees of freedom.

\begin{remark} \label{remark:convexifying_term}
	Although an approximation, the QoN cost is easily handled in the proposed method as a convex function of the variables (also by the block matrix methods \cite{ridderhofChanceConstrainedCovarianceControl2020,oguriChanceConstrainedControlSafe2024a}). The seminal work \cite{liuOptimalCovarianceSteering2025} relies on an assumption on the objective function for convexification, and we are not aware of a convex upper bound for this formulation. Workarounds such as linearizarion break the proof of lossless convexification in \cite{liuOptimalCovarianceSteering2025}. To retain the proof, an extra regularization term \(\eta \trace(Y_k)\) with small \( \eta > 0 \) can be added, as we do in \cite{kumagaiRobustCislunarLowThrust2025} ($Y_k := \cov{u_k}$).
\end{remark}

\subsection{Constraints}
The square root representation allows all other constraints to be rewritten as convex functions of the variables.

The structure of \(S_k\) is imposed by
\begin{equation} \label{eq:S_structure_constraint}
	S_k \in \mathbb{L}^n, \quad \mathrm{diag}(S_k) \geq 0, \quad k \in [N].
\end{equation}
The initial distributional constraint is imposed by
\begin{equation} \label{eq:initial_distributional_constraint}
	\mu_0 = \mu_{\mathrm{init}}, \quad S_0 = P_{\mathrm{init}}^{1/2}.
\end{equation}
The terminal distributional constraint is imposed by \cite{okamotoOptimalCovarianceControl2018}
\begin{equation} \label{eq:terminal_distributional_constraint}
	\mu_N = \mu_{\mathrm{fin}}, \quad \|P_{\mathrm{fin}}^{-1/2} S_N\|_2 - 1 \leq 0.
\end{equation}
The affine CC in \cref{eq:linear_chance_constraint} can be exactly expressed as \cite{okamotoOptimalCovarianceControl2018}:
\begin{equation} \label{eq:linear_chance_constraint_deterministic}
	\alpha_j^\top \mu_k + \mathcal{Q}_{\mathcal{N}}(1 - p_j) \|S_k^\top \alpha_j\|_2 - \beta_j \leq 0, \quad k \in [N]
\end{equation}
where \( \mathcal{Q}_{\mathcal{N}}(\cdot) \) is the inverse of the standard normal cumulative distribution function.
For the norm CC, similar to the QoN cost, the CC can be conservatively imposed as \cite{oguriChanceConstrainedControlSafe2024a}:
\begin{equation} \label{eq:norm_chance_constraint_deterministic}
	\norm{\mu_k}_2 + \sqrt{Q_{\chi^2_n}(1 - p_\gamma)} \ \norm{S_k}_2 - \gamma \leq 0, \quad k \in [N]
\end{equation}
To impose control CCs, replace \( \mu_k \) with \( v_k \) and \( S_k \) with \( L_k \).

To summarize, \cref{eq:problem_statement} is converted into a nonconvex optimization problem as follows:
\begin{equation} \label{eq:nonconvex_problem_statement}
	\minimize_{\bm{v}, \bm{L}, \bm{\mu}, \bm{S}} \cref{eq:objective_deterministic} \quad \text{s.t.} 
	\quad 
	\text{\cref{eq:S_structure_constraint,eq:mean_dynamics,eq:square_root_covariance_propagation,eq:initial_distributional_constraint,eq:terminal_distributional_constraint,eq:linear_chance_constraint_deterministic,eq:norm_chance_constraint_deterministic}}
\end{equation}
where the only source of nonconvexity is the square root covariance dynamics \cref{eq:square_root_covariance_propagation}. For notational simplicity, we define \(\bm{L} := \{L_k\}_{k=0}^{N-1}\), \(\bm{\mu} := \{\mu_k\}_{k=0}^{N}\), and \(\bm{S} := \{S_k\}_{k=0}^{N}\).